\section{A regular hypergraph of affine lines}\label{sec:hypergraph}

The algebraic criterion is now complete. Its remaining input is a line
system satisfying Proposition~\ref{prop:pairing}: full lines avoiding a
set of holes of controlled size, with exact regularity and incidence
girth at least $12$.
Starting from line partitions gives regularity, but simply deleting the
lines through a hole would lower the degree of other points as well.
We instead arrange the lines through each hole in planar bundles that
can be replaced by parallel lines missing the hole. Each replacement
preserves every other point's degree. We choose the holes to make these
replacements compatible, then use random three-dimensional slabs to
exclude the remaining short incidence cycles.

Let $q=2^h$ and let $\ell$ be the least prime congruent to $-1$ modulo
$q^{100}$. The residue $q^{100}-1$ is coprime to $q^{100}$. Linnik's theorem~\cite{Linnik1944} gives absolute constants $c,L>0$ such that
this prime exists and $\ell\le c q^{100L}$; see, for example,
\cite[Theorem~1.1]{Xylouris2011}.  Set
\begin{equation}\label{eq:hypergraph-parameters}
 D=\ell+1,\qquad s=\lfloor\log_qD\rfloor,\qquad
 b=s+20,\qquad n=10b,\qquad N_{\rm pt}=q^n.
\end{equation}
Thus $q^{100}\mid D$, $s\ge100$, and
\begin{equation}\label{eq:hypergraph-growth}
 q^s\le D<q^{s+1}.
\end{equation}
Moreover, $s$ is bounded independently of $h$: indeed,
$D\le(c+1)q^{100L}$ implies
$s\le100L+\log_2(c+1)$.  In this section all asymptotic statements refer
to $h\to\infty$.  Constants implicit in $O(\cdot)$ may depend on a
fixed upper bound for $s$, but not on $h$, the field enumeration,
or the indices of the objects being counted. Thus $s$ and $n$ may vary
with $h$, but the estimates below use fixed upper bounds for them.

The divisibility $q\mid D$ permits bundles of $q$ directions. The lower
bound $s\ge100$ makes the hole-selection estimates small, and the slack
in $b=s+20$ controls the later event weights. The factor ten in $n=10b$
permits ten-wise independent subspaces in the construction below. Each
step of a short polygon will involve at most two of these subspaces,
so their directness will cover polygons with at most five steps.

Let $Q$ be a field of order $q$, with any enumeration
$(a_0,\ldots,a_{q-1})$, and put $\mathcal A=Q^n$.  Define
\begin{equation}\label{eq:hypergraph-lower-set}
 d_0=(q-2)/2,\qquad
 S=\left\{(a_{i_1},\ldots,a_{i_n}):
                  i_1+\cdots+i_n\le d_0\right\}.
\end{equation}
An affine line always means a full translate of a one-dimensional
$Q$-subspace, and consequently contains exactly $q$ points.

\begin{theorem}\label{thm:hypergraph}
For all sufficiently large $h$, for every $Q$ and enumeration as above,
there are a set $O\subset S$ and a set $\mathcal F$ of affine lines
contained in $V=\mathcal A\setminus O$ such that
\begin{equation}\label{eq:hypergraph-hole-size}
 \frac{2N_{\rm pt}}{D^2}\le |O|\le\frac{7N_{\rm pt}}{D^2},
\end{equation}
every point of $V$ lies on exactly $D$ members of $\mathcal F$, and
the bipartite incidence graph of $V$ and $\mathcal F$ has girth at
least $12$.
\end{theorem}

\subsection{Direction spaces and the local trade}

Divide an index set of size $D$ into $D/q$ bundles of size $q$.
Write $\rho(j)$ for the bundle containing an index $j$.

\begin{lemma}\label{lem:hypergraph-direction-spaces}
There are two-dimensional subspaces $W_r\subset\mathcal A$, one per
bundle, and $(b-1)$-dimensional subspaces $E_j\subset\mathcal A$,
one per index, such that any at most ten distinct members of this list
have direct sum.  Choose a direction $e_r$ in each $W_r$, and assign
its other $q$ directions bijectively to the indices in that bundle,
calling them $d_j$.  Then $U_j=d_j\oplus E_j$ has dimension $b$, and
\begin{equation}\label{eq:hypergraph-intersections}
 U_j\cap U_{j'}=0\quad(j\ne j'),\qquad
 U_j\cap e_r=0,\qquad
 U_j\cap W_r=
 \begin{cases}d_j,&r=\rho(j),\\0,&r\ne\rho(j).
 \end{cases}
\end{equation}
\end{lemma}

\begin{proof}
Let $Q'/Q$ be an extension of degree $b$, and identify the $Q$-space
$\mathcal A$ with $(Q')^{10}$.  For distinct $\xi\in Q'$, the vectors
$(1,\xi,\ldots,\xi^9)$ have the property that any at most ten are
linearly independent over $Q'$: the first $k$ coordinates of $k$ such
vectors form an invertible Vandermonde matrix.  There are enough
parameters because
$D+D/q<q^{s+2}\le q^b$.  Give each $W_r$ and $E_j$ its own parameter,
and choose it, of the prescribed $Q$-dimension, inside the corresponding
one-dimensional $Q'$-space.  These choices have the required directness.

To verify the intersections, any equality involving $U_j$ and $U_{j'}$
first has zero $E_j$- and $E_{j'}$-components.  Its remaining components
belong either to independent $W$-spaces or to distinct directions in
one $W$-space.  In both cases they vanish.  The other assertions follow
the same way; within $W_{\rho(j)}$ the directions $d_j$ and $e_{\rho(j)}$
are distinct.
\end{proof}

\Needspace{10\baselineskip}
Fix these spaces for the rest of the construction. The eventual
preliminary partition indexed by $j$ will use varying directions in
$U_j$, with $o+d_j$ its line through a hole $o$; the spaces $E_j$
provide the other directions. If these $q$ axes are present for a
bundle $r$, then inside $o+W_r$ they can be replaced by the $q-1$
lines parallel to $e_r$ other than $o+e_r$. Each point off $o+e_r$
loses and gains one incidence, the hole loses $q$, and the remaining
points of $o+e_r$ are unaffected. Figure~\ref{fig:hypergraph-trade}
illustrates this local count. The hole restrictions and slab construction
below will supply these axes and make all the trades simultaneously
possible.
\begin{figure}[htbp]
\centering
\begin{tikzpicture}[>=Stealth,scale=.9,
  every node/.style={font=\small}]
 \begin{scope}[shift={(-3.35,0)}]
  \draw[rounded corners,gray!55] (-2.25,-1.65) rectangle (2.25,1.65);
  \node at (0,2.05) {Remove a fan of $q$ lines};
  \draw[gray,dashed] (0,-1.45)--(0,1.45);
  \draw[blue!65!black] (-1.9,-1.1)--(1.9,1.1);
  \draw[blue!65!black] (-1.9,-.4)--(1.9,.4);
  \draw[blue!65!black] (-1.9,.4)--(1.9,-.4);
  \draw[blue!65!black] (-1.9,1.1)--(1.9,-1.1);
  \fill[red!75!black] (0,0) circle (2.5pt);
  \node[fill=white,inner sep=1pt,above left] at (-.05,.08) {$o$};
  \node[fill=white,inner sep=1pt,above right] at (0,.98) {$o+e_r$};
  \node at (0,-2.05) {All directions in $W_r$ except $e_r$};
 \end{scope}
 \draw[->,thick] (-.65,0)--(.65,0);
 \begin{scope}[shift={(3.35,0)}]
  \draw[rounded corners,gray!55] (-2.25,-1.65) rectangle (2.25,1.65);
  \node at (0,2.05) {Insert $q-1$ parallel lines};
  \draw[gray,dashed] (0,-1.45)--(0,1.45);
  \draw[teal!70!black] (-1.65,-1.4)--(-1.65,1.4);
  \draw[teal!70!black] (-.8,-1.4)--(-.8,1.4);
  \draw[teal!70!black] (1.35,-1.4)--(1.35,1.4);
  \fill[red!75!black] (0,0) circle (2.5pt);
  \node[fill=white,inner sep=1pt,above left] at (-.05,.08) {$o$};
  \node[fill=white,inner sep=1pt,above right] at (0,.98) {$o+e_r$};
  \node at (0,-2.05) {Direction $e_r$; omit the line through $o$};
 \end{scope}
\end{tikzpicture}
\caption{Schematic of the trade inside $o+W_r$.
Every point off $o+e_r$ loses and gains one incidence;
$o$ loses $q$ incidences, and other points of $o+e_r$ are unaffected.
The strokes illustrate line families, not a real-plane realization
of finite-field incidences.}
\label{fig:hypergraph-trade}
\end{figure}


\subsection{Choosing compatible holes}

The replacement lines are parallel within a bundle, but lines from two
bundles could form a rectangle. We prevent this while choosing the holes.
An \emph{exceptional rectangle} has corners
$x,x+u,x+u+v,x+v$, where $0\ne u\in e_r$, $0\ne v\in e_{r'}$,
and $r\ne r'$.  For each of its four side lines, a \emph{witness}
is a point off that line but in the translate of $W_r$ or $W_{r'}$
containing it, according to the side's direction.
If that side is produced by a trade, its hole is such a witness.

\Needspace{13\baselineskip}
\begin{lemma}\label{lem:hypergraph-holes}
For all sufficiently large $h$, there is $O\subset S$ satisfying
Equation~\eqref{eq:hypergraph-hole-size} and the following properties.
\begin{enumerate}[label=\textup{(\roman*)}]
 \item Every translate of every $W_r$ contains at most one point of $O$.
 \item For every $j$, every coset $C$ of $U_j$, and every $r$,
       $|O\cap(C+W_r)|\le100$.
 \item No exceptional rectangle has four distinct witnesses in $O$,
       one for each of its four side lines.
\end{enumerate}
\end{lemma}

\begin{proof}
Write $M=N_{\rm pt}/D^2$.  Since $d_0<q$, counting nonnegative index
tuples gives $|S|=\binom{d_0+n}{n}$.  For large $h$, $d_0\ge q/3$,
so
\[
 |S|\ge \frac{q^n}{3^n n!}\ge c_S N_{\rm pt}
\]
with a constant $c_S>0$ independent of $h$, because $n$ is bounded.
Choose each point of $S$ independently with probability
$\theta=5M/|S|$.  For large $h$ this is at most $1$, and
$\theta=O(D^{-2})$.  If $X$ is the chosen set, then
$\mathbb E|X|=5M$ and $\operatorname{Var}(|X|)\le5M$.
Also $M>q^{n-2s-2}=q^{8s+198}\to\infty$.
Chebyshev's inequality therefore gives
\begin{equation}\label{eq:hypergraph-initial-holes}
 \Pr(4M\le |X|\le6M)=1-o(1).
\end{equation}

There are at most $D^2N_{\rm pt}$ triples $(j,r,C)$, and every
$C+W_r$ has at most $q^{b+2}$ points.  Repeated affine spaces in
this count only cause overcounting.  A union bound over choices of
$100$ selected points bounds the probability that any such space
contains more than $100$ points of $X$ by
\[
 D^2N_{\rm pt}(\theta q^{b+2})^{100}
 =O\left(q^{2s+2+n+100(b+2-2s)}\right)
 =O\left(q^{2402-88s}\right)=o(1).
\]

Mark every selected point belonging to a pair of distinct selected
points in one translate of some $W_r$.  Also mark every selected
point belonging to a tuple of four distinct selected witnesses for
an exceptional rectangle.  The expected number of collision pairs
is at most
$N_{\rm pt}(D/q)q^2\theta^2$.  There are at most
$N_{\rm pt}(D/q)^2(q-1)^2\le N_{\rm pt}D^2$
choices of rectangle data, and at most $q^8$ witness tuples per
rectangle.  Each distinct tuple is selected with probability at most
$\theta^4$; tuples containing a point outside $S$ have probability zero.
Thus the number $Z$ of marked points satisfies
\[
 \frac{\mathbb EZ}{M}
 \le \frac{2N_{\rm pt}(D/q)q^2\theta^2
             +4N_{\rm pt}D^2q^8\theta^4}{M}
 =O\left(\frac qD+\frac{q^8}{D^4}\right)=o(1).
\]
Markov's inequality gives $Z\le M$ with probability $1-o(1)$.
Choose an outcome satisfying this bound, Equation~\eqref{eq:hypergraph-initial-holes},
and the occupancy bound, and delete all marked points.  The remaining
set $O$ has $3M\le |O|\le6M$, hence the required bounds.  Deletion
preserves the occupancy bound and eliminates every obstruction in
\textup{(i)} and \textup{(iii)}.
\end{proof}

Fix such an $O$.  Put
\[
 T_r=\bigcup_{o\in O}(o+W_r).
\]
The \emph{replacement lines of label $R_r$} are all the lines of
direction $e_r$ in each $o+W_r$, except $o+e_r$ itself.

\begin{lemma}\label{lem:hypergraph-replacement-lines}
The replacement lines avoid $O$ and are pairwise disjoint within
each label $R_r$.  No exceptional rectangle has all four side lines
among these replacement lines.
\end{lemma}

\begin{proof}
Property~\textup{(i)} of Lemma~\ref{lem:hypergraph-holes} says that
different holes lie in disjoint translates of a fixed $W_r$.
Within one such plane the replacement lines are parallel and omit
its unique hole.  This proves the first assertions.

If a rectangle used four replacement lines, each side would have
its corresponding hole as a witness.  These four witnesses must
be distinct.  Opposite sides of direction $e_r$ lie in $W_r$-cosets
separated by a nonzero vector in $e_{r'}$, which lies outside $W_r$.
Their planes are therefore disjoint.  The planes of two adjacent
sides intersect exactly at their shared corner, because
$W_r\cap W_{r'}=0$.  This corner is on both side lines, so neither
off-line witness can equal it.  The resulting four distinct witnesses
contradict property~\textup{(iii)} of Lemma~\ref{lem:hypergraph-holes}.
\end{proof}

\subsection{Random slabs and short polygons}

The holes and the replacement lines are now fixed. It remains to
partition the rest of the space into lines without introducing short
polygons. We first choose the slabs in which those lines will lie,
excluding a finite list of events that would allow such polygons;
the line partitions themselves are constructed only afterward.

Independently for each index $j$ and each coset $C$ of $U_j$, choose
a uniformly random two-dimensional subspace $P_{j,C}\subset E_j$.
The translates of $d_j\oplus P_{j,C}$ partition $C$ into
three-dimensional affine \emph{slabs}.  For $u\in U_j$, write
$u^{(E_j)}$ for its
$E_j$-component.  A fixed nonzero vector of $E_j$ belongs to $P_{j,C}$
with probability
\begin{equation}\label{eq:hypergraph-inclusion-probability}
 p_* = \frac{q^2-1}{q^{b-1}-1}\le 2q^{3-b}.
\end{equation}
Indeed, linear automorphisms of $E_j$ act transitively on its nonzero
vectors, and every two-dimensional subspace contains $q^2-1$ of them.

We define a finite list of bad events before selecting any lines
inside the slabs.  First, for each pair of holes in a coset $C$ of
$U_j$, declare it bad if they occupy the same slab.  Their difference
has nonzero $E_j$-component: otherwise it would lie in
$d_j\subset W_{\rho(j)}$, contrary to
Lemma~\ref{lem:hypergraph-holes}\textup{(i)}.  This event is exactly
the inclusion of that component in $P_{j,C}$ and has probability $p_*$.

The other events concern closed labelled sequences of $k$ nonzero
steps, $k\in\{3,4,5\}$, in $\mathcal A$.  The allowed step types are
\begin{align*}
 G_j &: \quad u\in U_j,\quad u^{(E_j)}\ne0;\\
 R_r &: \quad u\in e_r\setminus\{0\},
               \quad\hbox{both endpoints belong to }T_r.
\end{align*}
We call $G_j$ steps \emph{generic} and $R_r$ steps \emph{replacement}
steps; a replacement step vector will also be called exceptional.
An $R_r$ step need not lie on a replacement line: its endpoint condition
only requires membership in $T_r$, which every actual replacement line
satisfies. This deliberately enlarges the list of possible polygons.
Require at least one $G$ step and distinct labels at cyclically
adjacent steps.  Repeated vertices are allowed, so this list can
only be larger than the list of simple polygons that will matter
for girth.  For each such sequence, declare it bad if every $G_j$
step has $u^{(E_j)}\in P_{j,C}$, where $C$ is the $U_j$-coset
containing that step.  All point, label, and step choices are finite,
so the complete event list is finite.

The three hole conditions now have separate roles. Condition~\textup{(i)}
separates the replacement planes within each bundle and makes each
hole-pair event test a nonzero $E_j$-component. Condition~\textup{(ii)}
controls the hole contributions to the event counts for a fixed
$P_{j,C}$. Condition~\textup{(iii)} has already excluded rectangles
formed entirely from replacement lines.

\Needspace{8\baselineskip}
\begin{lemma}\label{lem:hypergraph-polygons}
Every candidate sequence with a generic step has one of the following
forms.
\begin{enumerate}[label=\textup{(\roman*)}]
 \item A rectangle alternating $G_j,R_r$, with opposite vectors paired.
 \item A rectangle alternating distinct $G_j,G_{j'}$, with opposite
       vectors paired.
 \item A pentagon containing two $G_j$ steps, two $G_{j'}$ steps,
       and one $R_r$ step, where $\rho(j)=\rho(j')=r$.
\end{enumerate}
If $g$ is the number of distinct generic labels, its bad event uses
$2g$ distinct random variables and has probability $p_*^{2g}$.
In case~\textup{(iii)}, the exceptional vector and one vector of each
generic label determine the other two vectors, once their order is fixed.
\end{lemma}

\begin{proof}
Before classifying the sequence, observe that each step involves
at most two members of the list of $W$- and $E$-spaces.  Thus at most
$2k\le10$ distinct spaces occur, and their directness is available
from Lemma~\ref{lem:hypergraph-direction-spaces}.
Closure forces the sum of the $E_j$-components to be zero for each
generic label.  Each such label therefore occurs at least twice.
Cyclic non-adjacency in a sequence of at most five steps allows
any label at most twice.  Consequently $g$ is $1$ or $2$, and each
generic label occurs exactly twice.

If $g=1$, the direction $d_j$ and the distinct occurring directions
$e_r$ are linearly independent.  In its own bundle this uses just the
two distinct directions $d_j,e_{\rho(j)}$; other bundles contribute
independent $W$-spaces.  Every exceptional label must therefore occur
twice as well.  Only case~\textup{(i)} is possible, and closure gives
the claimed opposite vectors.

If $g=2$ and $k=4$, the labels must alternate, and
$U_j\cap U_{j'}=0$ forces the opposite vectors to pair.
If $k=5$, there is one exceptional step, of label $R_r$.
Its nonzero $e_r$-component must be canceled by the $W$-components
of the generic steps.  Neither zero contribution nor a contribution
from just one of the directions $d_j,d_{j'}$ can cancel it, since
each such direction differs from $e_r$.  Thus both generic labels belong
to bundle $r$, giving case~\textup{(iii)}.  After choosing the exceptional
vector and one vector of each generic label, closure prescribes the
sum of the remaining vector in $U_j$ and the remaining vector in
$U_{j'}$.  Their direct sum makes these vectors unique.

Finally, between the two occurrences of a generic label $G_j$ there
is, in one cyclic direction, exactly one intervening step.  That
step has a different label and a nonzero vector outside $U_j$, by
Equation~\eqref{eq:hypergraph-intersections}.  The two occurrences
therefore belong to different $U_j$-cosets: if the first starts at $x$
with vector $u\in U_j$ and the intervening vector is $v\notin U_j$,
the second starts at $x+u+v\notin x+U_j$. Variables from different
generic labels are also distinct.  All $2g$ variables are consequently
independent, and Equation~\eqref{eq:hypergraph-inclusion-probability}
gives the asserted event probability.
\end{proof}

\subsection{Avoiding the slab events}

The classification controls the probability of each polygon event. To
avoid all events simultaneously, we must also control how many of them
use any one slab variable. We bound the sum of their probabilities at
a fixed $P_{j,C}$; the local lemma will apply once this bound is at most
$1/16$.

\begin{lemma}\label{lem:hypergraph-event-bounds}
Let $\mu$ be the largest, over all variables $P_{j,C}$, of the sum
of probabilities of bad events using that variable.  Then
$\mu=O(q^{-7})$, uniformly in the variable, and every event uses
at most four variables.
\end{lemma}

\begin{proof}
The assertion about four variables follows from
Lemma~\ref{lem:hypergraph-polygons}, with one variable for a hole-pair
event.  Fix $P_{j,C}$.  There are at most $100$ holes in $C$, since
$C\subset C+W_r$ for any $r$.  Hole pairs therefore contribute at most
$100^2p_*$.

For a polygon using this variable, cyclically rotate so that one of
its $G_j$ steps starts the sequence.  Rotations and orders of at most
five labels introduce only an absolute constant factor in all the
following counts.  For a generic rectangle choose the other label
($\le D$ choices), the sequence's initial point in $C$ ($q^b$ choices), and
one vector of each label ($\le q^{2b}$ choices).  Closure determines
the other two vectors.  The contribution is therefore
$O(Dq^{3b}p_*^4)$.

For a mixed rectangle $G_j,R_r$, both endpoints of the initial
generic step belong to $C\cap T_r$.  A hole $o$ can contribute to
this intersection only if $o\in C+W_r$, so there are at most $100$
such holes.  Each intersection $C\cap(o+W_r)$, when nonempty, is
a coset of $U_j\cap W_r$.  It has $q$ points if $r=\rho(j)$ and one
point otherwise. Thus $C\cap T_r$ has at most $100q$ points in the
first case and at most $100$ in the second. Choose the two endpoints,
and then choose the
exceptional vector in $e_r$ in at most $q$ ways.  Summing over the
$D/q$ bundles gives at most
\[
 O\left(100^2\left(q^3+\frac Dq\,q\right)\right)
   =O\bigl(100^2(q^3+D)\bigr)
\]
mixed rectangles, and hence a contribution
$O(100^2(q^3+D)p_*^2)$.  This is an upper bound even when some
selected endpoints or vectors fail the candidate conditions.

For a pentagon the other generic label is in bundle $\rho(j)$,
giving at most $q$ choices.  Choose the sequence's initial point in $C$,
one vector of each generic label, and the exceptional vector.
There are at most $q^{b+2b+1}$ choices after the label is chosen.
Lemma~\ref{lem:hypergraph-polygons} determines the remaining vectors,
so the contribution is $O(q^{3b+2}p_*^4)$.

In total,
\begin{align*}
 \mu
 &\le 100^2p_*+
 O\left(Dq^{3b}p_*^4+(D+q^3)p_*^2+q^{3b+2}p_*^4\right)\\
 &=O\left(q^{-s-17}+q^{-7}+q^{-s-33}
                      +q^{-2s-31}+q^{-s-6}\right)
 =O(q^{-7}),
\end{align*}
using $b=s+20$, Equation~\eqref{eq:hypergraph-growth}, and
Equation~\eqref{eq:hypergraph-inclusion-probability}.  The constants
in these event counts are absolute; in particular the bound is uniform.
\end{proof}

The following is a special case of the asymmetric Lov\'asz local
lemma; we include the classical conditional-probability proof.
For the origins see Erd\H{o}s and Lov\'asz~\cite{ErdosLovasz1975};
for the product-form condition used below see
\cite[Theorem~1.1]{MoserTardos2009}. No resampling algorithm is needed.

\begin{lemma}\label{lem:hypergraph-local-lemma}
Let a finite collection of events be determined by independent
random variables, with each event assigned a nonempty set of at most
four variables that determine it.
If the sum of event probabilities at every variable is at most
$1/16$, there is positive probability that no event occurs.
\end{lemma}

\begin{proof}
Connect two events when their assigned sets share a variable.  An event
is independent of all its non-neighbors jointly.  Set $x_A=2\Pr(A)$.
Since each event has a nonempty assigned set, $0\le x_A\le1/8<1$,
and for every event
\[
 \sum_{B\text{ neighbor of }A}x_B\le 2\cdot4\cdot\frac1{16}
 =\frac12.
\]
It follows that
\begin{equation}\label{eq:hypergraph-local-lemma-condition}
 \Pr(A)\le x_A\prod_{B\text{ neighbor of }A}(1-x_B),
\end{equation}
because the product is at least $1-\sum_Bx_B\ge1/2$.

We prove inductively that
\[
 \Pr\left(A\,\middle|\,\bigcap_{B\in\mathcal S}B^c\right)\le x_A
\]
whenever $\mathcal S$ is a sublist excluding $A$, and that all
conditioning events in this assertion have positive probability.
The empty sublist follows from
Equation~\eqref{eq:hypergraph-local-lemma-condition}.
At each induction step, positivity for a sublist follows by exposing
its events in order and using the smaller conditional bounds and
$x_B<1$.  For the bound, split $\mathcal S$ into neighbors and
non-neighbors of $A$.  Let $B_0$ be avoidance of the non-neighbors
and $D_0$ avoidance of the neighbors.  Independence gives
$\Pr(A\mid B_0)=\Pr(A)$, and
\[
 \Pr(A\mid B_0\cap D_0)
 \le \frac{\Pr(A)}{\Pr(D_0\mid B_0)}
 \le \frac{\Pr(A)}
 {\prod_{B\in\mathcal S,\ B\text{ neighbor of }A}(1-x_B)}
 \le x_A.
\]
The middle bound uses the smaller conditional bounds successively;
the last uses Equation~\eqref{eq:hypergraph-local-lemma-condition}.
The same successive conditioning for the full finite event list
gives positive probability of avoiding all events.
\end{proof}

By Lemmas~\ref{lem:hypergraph-event-bounds} and
\ref{lem:hypergraph-local-lemma}, for all sufficiently large $h$ we
can fix the spaces $P_{j,C}$ so that every slab contains at most one
hole and none of the candidate polygon events occurs.

\subsection{Line partitions, the local trade, and girth}

The slab choices have excluded every candidate short polygon with a
generic step, but no line partition has yet been specified. We now
partition each slab, perform the trades that remove the holes, and
check that every possible short incidence cycle is one of the
excluded polygons.

\begin{lemma}\label{lem:hypergraph-slab-partition}
Let a three-dimensional affine space have direction $P\oplus d$,
where $\dim_QP=2$ and $\dim_Qd=1$.  If it has no designated point,
it has a line partition whose directions project nontrivially to $P$.
If it has a designated point $o$, it has a line partition containing
$o+d$, with every other line's direction projecting nontrivially to $P$.
\end{lemma}

\begin{proof}
In the first case use parallel lines of any fixed nonzero direction
in $P$.  In the second choose $0\ne\delta\in d$ and a $Q$-linear
operator $J$ on $P$ with no eigenvector over $Q$.  Identify $P$ with
the $Q$-space of a quadratic field extension of $Q$, and let $J$ be
multiplication by an element $\eta$ outside $Q$.  An equation
$\eta x=\lambda x$ with $x\ne0$ and $\lambda\in Q$ would imply
$\eta=\lambda$, so this operator has the required property.
Besides $o+d$, use the lines
\begin{equation}\label{eq:hypergraph-slab-lines}
 \left\{o+x+\gamma(Jx)\delta:
                   x\in P,\ \gamma(x)=1\right\},
 \qquad 0\ne\gamma\in\operatorname{Hom}_Q(P,Q).
\end{equation}
The set $\gamma(x)=1$ is an affine line in $P$.  Its displayed lift
is injective, with direction projecting onto the nonzero line
$\ker\gamma$, so it is an affine line with the required direction.
For each $x\ne0$, the vectors $x,Jx$ form a basis of $P$.
For every $c\in Q$ there is consequently exactly one functional
with $\gamma(x)=1$ and $\gamma(Jx)=c$, and it is nonzero.
Thus the $q^2-1$ lines in Equation~\eqref{eq:hypergraph-slab-lines}
partition the $q^3-q$ points off $o+d$.  The axis $o+d$ supplies the
remaining $q$ points.
\end{proof}

For each $j$, apply Lemma~\ref{lem:hypergraph-slab-partition}
separately in every slab of every coset of $U_j$, using its unique
hole as the designated point when present.  This gives a preliminary
line partition of all of $\mathcal A$ for each $j$.  The line through
a hole $o$ in partition $j$ is $o+d_j$.  These hole axes are the only
lines whose directions have zero $E_j$-component.  Every other line
has a nonzero $E_j$-component in its direction, lying in the
corresponding $P_{j,C}$.

\begin{lemma}\label{lem:hypergraph-trade}
For every hole $o$ and bundle $r$, remove the $q$ lines $o+d_j$ with
$\rho(j)=r$, and insert the $q-1$ replacement lines of label $R_r$
in $o+W_r$.  After all these operations, the resulting set
$\mathcal F$ consists of lines in $V$, and every point of $V$ has
degree $D$.  The surviving lines of each preliminary partition are
pairwise disjoint, as are the replacement lines within each label.
There are no duplicate lines across different labels.
\end{lemma}

\begin{proof}
Within $o+W_r$, the removed fan uses every direction except $e_r$.
Each point off $o+e_r$ lies on exactly one removed line and exactly
one inserted line.  The hole $o$ loses $q$ incidences and gains none;
other points of $o+e_r$ lose and gain none.  Thus every non-hole point
has zero net change in degree in each trade.

No removal is requested twice.  Two holes using the same line of
one preliminary partition would be in the same slab, which was
excluded.  Lines from distinct preliminary partitions have directions
in distinct $U_j$'s with zero intersection, so cannot coincide.
Inserted lines of a given label are distinct by
Lemma~\ref{lem:hypergraph-replacement-lines}.  Lines of different
labels cannot coincide either: their direction spaces are distinct
$U_j$'s or $e_r$'s with pairwise zero intersections, by
Lemma~\ref{lem:hypergraph-direction-spaces}.
Hence the pointwise degree changes genuinely add when all trades
are performed, even if planes from different bundles overlap.
Initially every point has one incidence in each of the $D$ partitions.
Each hole loses all $D$ of its incidences, and every non-hole retains
degree $D$.  All inserted lines avoid the holes, and all preliminary
lines through a hole were removed.  Finally, disjointness within a
surviving preliminary label is inherited from its partition, and
disjointness within a replacement label was proved in
Lemma~\ref{lem:hypergraph-replacement-lines}.
\end{proof}

The pointwise incidence changes are summarized in
Figure~\ref{fig:hypergraph-trade}.

\begin{proof}[Proof of Theorem~\ref{thm:hypergraph}]
Use the holes from Lemma~\ref{lem:hypergraph-holes}, the successful
slab choices, and the line family from Lemma~\ref{lem:hypergraph-trade}.
The size and degree assertions have been proved.  Label surviving
preliminary lines from partition $j$ by $G_j$, and replacement lines
by their labels $R_r$.

Two distinct affine lines have at most one common point, so a cycle
of length less than $12$ in the incidence graph would give a polygon
of $k=3,4$, or $5$ nonzero steps along its successive lines.
Consecutive labels, including the last and first, are different
because lines of one label are disjoint.  If there is a generic step,
this polygon is one of the candidate sequences used above.  Each
generic step's $E_j$-component lies in its slab direction $P_{j,C}$,
so the corresponding bad event would occur, a contradiction.

If all steps have replacement labels, the distinct occurring $e_r$
directions are independent: at most five $W_r$'s occur.
Closure forces every label to occur at least twice, while cyclic
non-adjacency allows at most twice.  The only possibility is a
four-step alternating rectangle, with opposite vectors paired.
This is forbidden by Lemma~\ref{lem:hypergraph-replacement-lines}.
Thus the incidence girth is at least $12$.
\end{proof}

This supplies the full-line system, exact degree, and girth required by
Proposition~\ref{prop:pairing}. That proposition uses the same holes to
produce the two matrix pairings and the size inequalities for
Theorem~\ref{thm:criterion}. Section~\ref{sec:prescription} now makes the
choices by finite enumeration.
